\documentclass[11pt]{beamer}
\usepackage[utf8]{inputenc}
\usepackage{amsmath, amssymb, graphicx}

% Presentation Theme
\usetheme{Madrid}
\usecolortheme{default}

\title[Interplanetary Trajectories II]{Astrodynamics: Space Engineering Vehicle and Mission Design}
\subtitle{Time of Flight, Hyperbolic Orbits, \& Planetary Encounters}
\author{Module: Astrodynamics}
\date{Week 5}

\begin{document}

% --- Title Slide ---
\begin{frame}
    \titlepage
\end{frame}

% --- Slide 1 ---
\begin{frame}{Recap: The Patched Conic Approximation}
    \textbf{The 3-Phase Interplanetary Journey:}
    \begin{enumerate}
        \item \textbf{Planetocentric Departure:} Spacecraft leaves Planet 1 on an \textit{escape hyperbola}.
        \item \textbf{Heliocentric Transfer:} Spacecraft travels between planetary orbits on an \textit{elliptical} trajectory governed by the Sun.
        \item \textbf{Planetocentric Arrival:} Spacecraft enters the Sphere of Influence (SOI) of Planet 2 on an \textit{arrival hyperbola}.
    \end{enumerate}
    
    \vspace{0.3cm}
    \textit{Today, we will focus on the exact timing required for this transfer and the geometry of hyperbolic encounters.}
\end{frame}

% --- Slide 2 ---
\begin{frame}{Heliocentric Transfer: Time of Flight \& Phasing}
    To reach another planet, the target must be in the exact right position when the spacecraft arrives. 
    
    \vspace{0.2cm}
    \textbf{1. Time of Flight (TOF):} \\
    For a Hohmann transfer, the spacecraft travels exactly half of an ellipse ($180^\circ$ or $\pi$ radians).
    \begin{equation}
        TOF = \pi \sqrt{\frac{a_{tx}^3}{\mu_{Sun}}}
    \end{equation}

    \textbf{2. Initial Phase Angle ($\Delta \theta_0$):} \\
    Planet 2 moves while the spacecraft is in transit. The required "head start" angle for Planet 2 at the moment of launch is:
    \begin{equation}
        \Delta \theta_0 = \pi - \omega_2 (TOF)
    \end{equation}
    
    % INSTRUCTOR NOTE: Insert your "Oportunidades de encuentro" image here.
    \begin{center}
        % \includegraphics[height=3cm]{phasing_angle_image.png} 
        \textit{[Insert Planetary Phasing Diagram Here]}
    \end{center}
\end{frame}

% --- Slide 3 ---
\begin{frame}{Phase 1: The Departure Hyperbola}
    Once the heliocentric transfer is designed, we calculate the required \textbf{Hyperbolic Excess Velocity ($v_\infty$)}. This is the speed the spacecraft must have \textit{after} escaping Earth's gravity.
    
    \begin{equation}
        v_{\infty} = | V_{tx} - V_{Earth} |
    \end{equation}
    
    \textbf{The Burnout Velocity ($v_p$):} \\
    To leave the parking orbit ($r_p$) and achieve $v_\infty$ at infinity, the engine must provide a $\Delta V$ to reach the periapsis velocity of the escape hyperbola:
    \begin{equation}
        v_p = \sqrt{v_\infty^2 + \frac{2\mu_1}{r_p}}
    \end{equation}

    % INSTRUCTOR NOTE: Insert your Earth Departure vector image here.
    \begin{center}
        % \includegraphics[height=3cm]{departure_hyperbola.png}
        \textit{[Insert Departure Hyperbola Diagram Here]}
    \end{center}
\end{frame}

% --- Slide 4 ---
\begin{frame}{Phase 3: The Arrival Hyperbola}
    The spacecraft arrives at Planet 2 with an arrival excess velocity:
    \begin{equation}
        v_{\infty, arr} = | V_{arr} - V_{Target} |
    \end{equation}
    
    As it crosses the Sphere of Influence and falls toward the planet, gravity accelerates the spacecraft. At the periapsis ($r_{cap}$), the velocity will naturally be:
    \begin{equation}
        v_{p, arr} = \sqrt{v_{\infty, arr}^2 + \frac{2\mu_2}{r_{cap}}}
    \end{equation}
    
    To avoid performing a flyby and flying back out into deep space, the spacecraft must fire its engines in reverse (retro-burn) to slow down to the circular capture velocity ($v_c = \sqrt{\mu_2/r_{cap}}$).
\end{frame}

% --- Slide 5 ---
\begin{frame}{Planetary Encounter (The Gravity Assist)}
    What if we \textit{don't} want to stay at Planet 2? We can use its gravity to change our heliocentric velocity for free!
    
    \vspace{0.2cm}
    \textbf{The Flyby Physics:}
    \begin{itemize}
        \item Relative to the planet, the incoming speed and outgoing speed are identical ($|v_{\infty, in}| = |v_{\infty, out}|$). Energy is conserved.
        \item However, the planet's gravity \textbf{bends the trajectory} by a turning angle ($\delta$).
        \item Relative to the \textbf{Sun}, this change in vector direction results in a massive gain (or loss) of absolute velocity.
    \end{itemize}

    % INSTRUCTOR NOTE: Insert your "Encuentros planetarios/Flybys" image here.
    \begin{center}
        % \includegraphics[height=3.5cm]{planetary_flyby.png}
        \textit{[Insert Planetary Encounter Geometry Here]}
    \end{center}
\end{frame}

% --- Slide 6 ---
\begin{frame}{Gravity Assist Mathematics}
    The effectiveness of a planetary encounter depends on how close the spacecraft dares to fly to the planet ($r_p$).
    
    \vspace{0.3cm}
    \textbf{The Turning Angle ($\delta$):} \\
    The angle by which the excess velocity vector is rotated depends on the planet's gravity ($\mu$), the flyby altitude ($r_p$), and the arrival speed ($v_\infty$).
    \begin{equation}
        \sin \left( \frac{\delta}{2} \right) = \frac{1}{1 + \frac{r_p v_\infty^2}{\mu}}
    \end{equation}

    \vspace{0.3cm}
    \textbf{Mission Design Insight:} \\
    To get the maximum "kick" (largest $\delta$) from a gravity assist, you want to fly as close to the planet as physically possible without hitting the atmosphere, and ideally pass by a very massive planet (like Jupiter).
\end{frame}

\end{document}